\section{Ορισμός του Προβλήματος}
\label{sec:intro_problem}

Το πρόβλημα που καλείται να επιλύσει η παρούσα διπλωματική εργασία διατυπώνεται τυπικά ως πρόβλημα μέγιστης βελτιστοποίησης μιας συνάρτησης μέλανης κουτιού (black-box):
\begin{equation}
\label{eq:opt_problem}
\max_{\mathbf{x} \in \mathcal{X}} F(\mathbf{x}) = \left(\frac{L}{D}\right)(\mathbf{x}) - \mathcal{P}_{\text{stab}}(\mathbf{x}) - \mathcal{P}_{\text{vol}}(\mathbf{x})
\end{equation}
όπου $\mathcal{X} \subset \mathbb{R}^{10}$ είναι ο εφικτός χώρος σχεδιασμού και $F(\mathbf{x})$ είναι η σύνθετη συνάρτηση καταλληλότητας που μεγιστοποιείται.

Το διάνυσμα γονιδίων ορίζεται ως:
\begin{equation}
\label{eq:gene_vector}
\mathbf{x} = \bigl(b,\; c_{\text{root}},\; c_{\text{tip}},\; \Lambda,\; \Gamma,\; \theta,\; M,\; P,\; T,\; L_{\text{fuse}}\bigr) \in \mathbb{R}^{10}
\end{equation}
όπου $b$ είναι το εκπέτασμα πτέρυγας, $c_{\text{root}}$ και $c_{\text{tip}}$ οι χορδές ρίζας και ακροπτερυγίου αντίστοιχα, $\Lambda$ η γωνία βέλους (sweep angle), $\Gamma$ το δίεδρο, $\theta$ η συστροφή (twist), $M$, $P$, $T$ οι παράμετροι αεροτομής NACA (μέση καμπυλότητα, θέση μέγιστης καμπυλότητας, πάχος), και $L_{\text{fuse}}$ το μήκος ατράκτου~\cite{kulfan2008cst,abbott1959theory}.

Οι δεσμεύσεις του προβλήματος περιλαμβάνουν δύο ποινικούς όρους:
\begin{itemize}
    \item \textbf{Δέσμευση στατικής ευστάθειας:} $\mathcal{P}_{\text{stab}}(\mathbf{x}) > 0$ εφόσον $C_{m_\alpha}(\mathbf{x}) \geq 0$, δηλαδή η κλίση της ροπής βάθους ως προς τη γωνία πρόσπτωσης πρέπει να είναι αρνητική (διαμήκης ευστάθεια).
    \item \textbf{Δέσμευση εσωτερικού όγκου:} $\mathcal{P}_{\text{vol}}(\mathbf{x}) > 0$ εφόσον $V_{\text{fuse}}(\mathbf{x}) < V_{\min} = 20.000\ \text{mm}^3$, εξασφαλίζοντας επαρκή χώρο για την ενσωμάτωση ηλεκτρονικών εξαρτημάτων.
\end{itemize}

Η κεντρική υπολογιστική πρόκληση προκύπτει από τρία αλληλένδετα υποπροβλήματα:
\begin{enumerate}
    \item \textbf{Κόστος Αξιολόγησης:} Η συνάρτηση $F(\mathbf{x})$ είναι μη-αναλυτική, πολυτροπική και υπολογιστικά ακριβή· ένας εξελικτικός αλγόριθμος χρειάζεται τυπικά $10^3$--$10^5$ αξιολογήσεις για να συγκλίνει~\cite{jin2011surrogate}, καθιστώντας απαραίτητη τη χρήση μοντέλων surrogate που αντικαθιστούν τον πλήρη προσομοιωτή όταν αυτό είναι εφικτό.
    \item \textbf{Επιμονή Γνώσης:} Κάθε αξιολόγηση αποτελεί πολύτιμη πληροφορία. Ένα σύστημα που δεν αποθηκεύει και δεν επαναχρησιμοποιεί ήδη υπολογισμένα αποτελέσματα σπαταλά πόρους επαναϋπολογίζοντας το $F(\mathbf{x})$ για παρόμοιες γεωμετρίες. Η κατανεμημένη αποθήκευση μέσω δομής DHT~\cite{stoica2001chord} αντιμετωπίζει αυτό το κενό.
    \item \textbf{Παραλληλοποίηση:} Ο αλγόριθμος νήσων (island-model GA) επιτρέπει την ταυτόχρονη εξέλιξη πολλαπλών υποπληθυσμών σε ανεξάρτητους κόμβους~\cite{cantupaz2000parallel,alba1999survey}, μειώνοντας τον πραγματικό χρόνο (wall-clock time) σύγκλισης μέσω οριζόντιας κλιμάκωσης.
\end{enumerate}

Η αντιμετώπιση και των τριών υποπροβλημάτων απαιτεί μία ενοποιημένη αρχιτεκτονική 3 επιπέδων, η οποία περιγράφεται λεπτομερώς στο Κεφάλαιο~\ref{ch:architecture}.

